\chapter{Phương pháp Đề xuất}
\label{ch:method}

Chương này trình bày chi tiết và toàn diện về mặt toán học cũng như thuật toán của phương pháp nghiên cứu đề xuất: \textbf{Kiến trúc Kiểm chứng Neuro-Symbolic Neo Cú pháp Trừu tượng (AST-Anchored Neuro-Symbolic Verification Framework)}. Phương pháp được thiết kế theo triết lý phân tách nghiêm ngặt giữa bộ sinh đề xuất không tin cậy (LLM) và bộ khung kiểm chứng hình thức đáng tin cậy (AST, Trusted Symbolic Executor, Z3 Solver, và CPython Runtime).

\section{Tổng quan Kiến trúc Hệ thống}
\label{sec:architecture_overview}

Kiến trúc tổng thể của hệ thống bao gồm một chuỗi xử lý liên hoàn gồm 8 giai đoạn khép kín, được tổ chức thành hai luồng song song gặp nhau tại chốt kiểm soát chất lượng \textbf{Anchor Gate}, như được trực quan hóa trong Hình \ref{fig:system_architecture}.

\begin{figure}[H]
\centering
\begin{tikzpicture}[
    box/.style={draw, thick, fill=blue!5, rounded corners, align=center,
                minimum height=0.8cm, inner xsep=5pt, font=\small},
    trusted/.style={box, draw=red!80!black, very thick, fill=red!5, font=\small\bfseries},
    untrusted/.style={box, draw=orange!80!black, dashed, fill=yellow!10},
    side/.style={text width=3.7cm},
    gate/.style={diamond, aspect=2.1, draw=purple!80!black, thick, fill=purple!10,
                 align=center, font=\small\bfseries, inner sep=2pt},
    arrow/.style={-{Stealth}, thick},
    edge label/.style={font=\scriptsize, inner sep=2pt}
]

\node[box] (code) at (0,0) {Mã nguồn Python};
\node[trusted] (subset) at (0,-1.15) {1. Safe Subset Filter (QF-LIA)};
\node[trusted, side] (ast) at (-4.65,-2.5) {2. AST \\ Canonicalization \& SSA};
\node[untrusted, side] (llm) at (4.65,-2.5) {4. LLM Structured CoT};
\node[trusted, side] (symexec) at (-4.65,-3.9) {3. Trusted Symbolic \\ Executor};
\node[gate] (gate) at (0,-3.9) {5. Anchor Gate};
\node[box, side] (inv) at (-4.65,-5.55) {7. Loop Invariants ($k$-ind)};
\node[trusted] (twoquery) at (0,-5.55) {6. Two-Query Z3 Protocol};
\node[trusted, side] (replay) at (4.65,-5.55) {8. CPython Replay};
\node[untrusted, side] (feedback) at (4.65,-7.0) {Retrospective Feedback};
\node[box, fill=green!10, draw=green!60!black, very thick]
    (verdict) at (0,-8.15) {Claim Results \\ \& Roll-up Metrics};
\coordinate (replay-lane) at (2.25,-6.32);

\draw[arrow] (code.south) -- (subset.north);
\draw[arrow] (subset.west) -| (ast.north);
\draw[arrow] (subset.east) -| (llm.north);
\draw[arrow] (ast.south) -- (symexec.north);
\draw[arrow] (symexec.east) -- (gate.west);
\draw[arrow] (llm.south) |- (gate.east);
\draw[arrow] (gate.south) -- node[edge label, left=1mm] {Anchored Claims} (twoquery.north);
\draw[arrow] (inv.east) -- (twoquery.west);
\draw[arrow] (twoquery.east) -- node[edge label, above=4.5mm] {SAT (Model)} (replay.west);
\draw[arrow] (replay.south) -- node[edge label, right] {Crash Trace} (feedback.north);
\draw[arrow, dashed, orange!80!black] (feedback.east) -- (7,-7) |- (llm.east);
\draw[arrow] (twoquery.south) -- (verdict.north);
\draw[arrow] ([xshift=-1.2cm]replay.south) |- (replay-lane) |- (verdict.east);

\end{tikzpicture}
\caption{Kiến trúc tổng thể của AST-Anchored Neuro-Symbolic Verification Framework}
\label{fig:system_architecture}
\end{figure}

Nguyên tắc phân định ranh giới độ tin cậy được thể hiện rõ:
\begin{itemize}
    \item \textbf{Cơ sở Tính toán Tin cậy (Trusted Computing Base -- TCB):} Bao gồm Bộ phân tích cú pháp AST, Bộ lọc tập con, Trình thực thi ký hiệu SSA, Cổng neo Anchor Gate, Bộ giải Z3 SMT, và Trình thông dịch CPython. Toàn bộ logic ra quyết định kiểm chứng đều nằm hoàn toàn trong TCB.
    \item \textbf{Thành phần Không Tin cậy (Untrusted Generator):} Mô hình ngôn ngữ lớn (Gemini) chỉ có quyền hạn duy nhất là sinh ra các chuỗi văn bản gợi ý (reasoning trace, ứng viên bất biến, hoặc đề xuất sửa sai). Mọi gợi ý này bắt buộc phải chịu sự phán quyết của TCB trước khi được chấp nhận.
\end{itemize}

\section{Giai đoạn 1: Lọc Tập con An toàn (Safe Subset Filtering)}
\label{sec:safe_subset}

Để bảo đảm tính khả chứng và ngăn ngừa lỗi sụp đổ không mong muốn do các tính năng ngôn ngữ chưa được mô hình hóa, hệ thống triển khai bộ lọc \texttt{SafeSubsetFilter} đóng vai trò là cửa ngõ đầu vào.

\subsection{Đặc tả Whitelist Ngữ pháp QF-LIA}
Một chương trình Python được chấp nhận vào pipeline kiểm chứng nếu và chỉ nếu toàn bộ các cấu trúc cú pháp của nó nằm hoàn toàn trong danh mục cho phép (whitelist):
\begin{enumerate}
    \item \textbf{Câu lệnh (\texttt{ast.stmt}):} \texttt{FunctionDef}, \texttt{Assign}, \texttt{AugAssign} (\texttt{+=}, \texttt{-=}), \texttt{If}, \texttt{While}, \texttt{For} (chỉ cho phép dạng duyệt dải số nguyên \texttt{range(stop)} hoặc \texttt{range(start, stop)}), \texttt{Pass}, \texttt{Return}, và \texttt{Assert}.
    \item \textbf{Biểu thức (\texttt{ast.expr}):} \texttt{Constant} (chỉ số nguyên $\mathbb{Z}$), \texttt{Name} (tên biến), \texttt{UnaryOp} (\texttt{UAdd}, \texttt{USub}, \texttt{Not}), \texttt{BinOp} (\texttt{Add}, \texttt{Sub}, \texttt{Mult} -- với điều kiện ít nhất một toán hạng là hằng số nguyên tĩnh), \texttt{BoolOp} (\texttt{And}, \texttt{Or}), và \texttt{Compare} (\texttt{Eq}, \texttt{NotEq}, \texttt{Lt}, \texttt{LtE}, \texttt{Gt}, \texttt{GtE}).
\end{enumerate}

Bất kỳ chương trình nào chứa các cấu trúc ngoài whitelist như: số thực (\texttt{float}), chuỗi ký tự (\texttt{str}), danh sách (\texttt{list}), gọi hàm ngoại vi, xử lý ngoại lệ (\texttt{try-except}), hoặc phép nhân hai biến ký hiệu ($x \cdot y$) sẽ bị từ chối ngay lập tức tại cổng lọc và gán trạng thái \texttt{REJECTED\_BY\_SUBSET}, hoàn toàn không được đưa vào solver.

\subsection{Thuật toán Chuẩn hóa Mã nguồn}
Thuật toán kiểm tra và chuẩn hóa được mô tả chi tiết trong Thuật toán \ref{alg:subset_filter}.

\begin{algorithm}[H]
\SetAlgoLined
\KwIn{Mã nguồn chương trình Python $S_{\text{code}}$}
\KwOut{Đối tượng kết quả $R_{\text{subset}} = \langle \text{is\_valid}, \text{normalized\_ast}, \text{reason} \rangle$}
\BlankLine
\Begin{
    $T \leftarrow \text{ast.parse}(S_{\text{code}})$\;
}{
    \Return $\langle \text{False}, \text{None}, \text{``Cú pháp không hợp lệ''} \rangle$\;
}
\ForEach{\text{node} $\in \text{ast.walk}(T)$}{
    \If{$\text{type}(\text{node}) \notin \text{WHITELIST\_NODES}$}{
        \Return $\langle \text{False}, \text{None}, \text{f``Nút không được hỗ trợ: } \{\text{type}(\text{node})\} \text{''} \rangle$\;
    }
    \If{$\text{isinstance}(\text{node}, \text{ast.BinOp})$ \textbf{and} $\text{isinstance}(\text{node.op}, \text{ast.Mult})$}{
        \If{\textbf{not} $(\text{isinstance}(\text{node.left}, \text{ast.Constant}) \lor \text{isinstance}(\text{node.right}, \text{ast.Constant}))$}{
            \Return $\langle \text{False}, \text{None}, \text{``Phép nhân phi tuyến (hai biến ký hiệu)''} \rangle$\;
        }
    }
}
$T_{\text{norm}} \leftarrow \text{NormalizeAugAssign}(T)$ \tcp*{Biến đổi x += 1 thành x = x + 1}
\Return $\langle \text{True}, T_{\text{norm}}, \text{``Hợp lệ''} \rangle$\;
\caption{Thuật toán Lọc và Chuẩn hóa Tập con An toàn (Safe Subset Filter)}
\label{alg:subset_filter}
\end{algorithm}

\section{Giai đoạn 2: Định danh AST Canonical và Phân tích Phạm vi Biến}
\label{sec:ast_canonicalization}

\subsection{Cấp phát Định danh Nút Ổn định (\texttt{assign\_ids})}
Mỗi nút trong cây AST đã chuẩn hóa được gán một thuộc tính duy nhất \texttt{node\_id}. Thuật toán duyệt cây theo thứ tự trước (pre-order traversal). Với một nút cú pháp $v$ bắt đầu tại dòng $L$, định danh chính tắc được thiết lập theo quy tắc:
\begin{equation}
\text{node\_id}(v) = \begin{cases}
\text{type}(v) \text{@L} \cdot L & \text{nếu } v \text{ là nút duy nhất cùng kiểu trên dòng } L \\
\text{type}(v) \text{@L} \cdot L \text{\_} k & \text{nếu } v \text{ là nút thứ } k \ (k \ge 1) \text{ cùng kiểu trên dòng } L
\end{cases}
\end{equation}

Định danh này bảo đảm tính đơn ánh: hai câu lệnh khác nhau trong cùng một chương trình luôn sở hữu hai \texttt{NodeId} phân biệt, cho phép truy vết chính xác nguồn gốc của mọi nghĩa vụ kiểm chứng.

\subsection{Bảng Ký hiệu và Phạm vi Từ vựng (Lexical Scope Analysis)}
Hệ thống xây dựng một Bảng Phạm vi Ký hiệu (Lexical Scope Table) $\mathcal{T}_{\text{scope}}: \text{NodeId} \rightarrow 2^{\mathcal{V}}$. Với mỗi nút câu lệnh $v$, $\mathcal{T}_{\text{scope}}(v)$ lưu trữ tập hợp tất cả các tên biến đã được khai báo hoặc gán giá trị trên ít nhất một đường thực thi dẫn tới $v$. Nếu một biến $y$ xuất hiện trong một biểu thức tại $v$ mà $y \notin \mathcal{T}_{\text{scope}}(v)$, hệ thống xác định đây là một biến chưa khởi tạo hoặc một biến ảo giác do mô hình tự nghĩ ra.

\subsection{Trích xuất Khẳng định Tiên nghiệm từ Mã nguồn}
Cùng với các khẳng định do LLM sinh ra, hệ thống tự động trích xuất tất cả các câu lệnh \texttt{assert Condition} có sẵn trong mã nguồn:
\begin{equation}
\text{Claim}_{\text{assert}} = \langle \text{NodeId}(\text{assert\_node}), \text{lineno}, \text{Condition}, \text{Source} = \text{``source\_code''} \rangle
\end{equation}
Bộ khẳng định tiên nghiệm này được dùng để xây dựng Hồ chứa Khẳng định Dùng chung (\textit{Shared Claim Pool}), bảo đảm tính công bằng tuyệt đối khi so sánh với các baseline.

\section{Giai đoạn 3: Thực thi Ký hiệu Đáng tin cậy (Trusted Symbolic Execution)}
\label{sec:trusted_symbolic_execution}

Mô-đun \texttt{TrustedSymbolicExecutor} là trái tim toán học của TCB, chịu trách nhiệm chuyển hóa mã nguồn thành hệ phương trình ràng buộc ngữ nghĩa hoàn chỉnh.

\subsection{Quy tắc Chuyển đổi SSA và Môi trường Biến}
Trạng thái ký hiệu tại một điểm thực thi được đặc trưng bởi hàm phiên bản biến $\sigma: \mathcal{V} \rightarrow \mathbb{N}$ và ánh xạ giá trị ký hiệu $\mathcal{E}: \mathcal{V}_{\text{SSA}} \rightarrow \text{Expr}_{\text{SMT}}$.

Với mỗi câu lệnh cú pháp $v$, bộ thực thi sinh ra một \textit{Quan hệ Chuyển trạng thái Nút} (Node Transition Relation) $T_v(s, s')$ theo các quy tắc toán học:

\begin{enumerate}
    \item \textbf{Câu lệnh Gán Biến ($x = e$):}
    \begin{align}
    \text{Tăng chỉ số:} \quad & k_{\text{new}} = \sigma(x) + 1, \quad \sigma' = \sigma[x \mapsto k_{\text{new}}] \\
    \text{Ràng buộc chuyển tiếp:} \quad & T_v(s, s') \equiv \Big( x^{(k_{\text{new}})} = \text{Eval}(e, \sigma) \Big) \land \bigwedge_{y \in \mathcal{V} \setminus \{x\}} (y^{(\sigma'(y))} = y^{(\sigma(y))})
    \end{align}
    Mọi biến khác không bị thay đổi giữ nguyên giá trị qua phép đồng nhất.

    \item \textbf{Khối Rẽ nhánh Điều kiện (\texttt{if $B$ then $C_{\text{then}}$ else $C_{\text{else}}$}):}
    Giả sử trạng thái trước khi rẽ nhánh là $\sigma_0$.
    \begin{itemize}
        \item Thực thi nhánh Then dưới tiền đề $\text{Guard}_{\text{then}} \equiv \text{Eval}(B, \sigma_0)$, thu được trạng thái biến $\sigma_1$.
        \item Thực thi nhánh Else dưới tiền đề $\text{Guard}_{\text{else}} \equiv \neg \text{Eval}(B, \sigma_0)$, thu được trạng thái biến $\sigma_2$.
    \end{itemize}
    Tại điểm hội tụ sau khối lệnh \texttt{if}, với mỗi biến $z \in \mathcal{V}$ bị cập nhật trong ít nhất một nhánh, một phiên bản mới $z^{(\sigma'(z))}$ được sinh ra với điều kiện chuyển tiếp ITE:
    \begin{equation}
    z^{(\sigma'(z))} = \text{ite}\Big(\text{Eval}(B, \sigma_0), \; z^{(\sigma_1(z))}, \; z^{(\sigma_2(z))}\Big)
    \end{equation}

    \item \textbf{Khối Vòng lặp (\texttt{while $B$ do $C_{\text{body}}$}):}
    Đối với các vòng lặp hữu hạn hoặc khi thực hiện duỗi vòng lặp có chặn (Bounded Loop Unrolling với độ sâu $K$):
    \begin{equation}
    T_{\text{loop}} \equiv \bigvee_{j=0}^{K} \left( \bigwedge_{i=0}^{j-1} \Big( \text{Guard}(s_i) \land T_{\text{body}}(s_i, s_{i+1}) \Big) \land \neg \text{Guard}(s_j) \right)
    \end{equation}
    Nếu sau $K$ bước lặp mà điều kiện dừng $\neg \text{Guard}$ chưa đạt được, trạng thái được đánh dấu là không đầy đủ (incomplete coverage), dẫn tới việc hạn chế kết luận \texttt{VERIFIED}.
\end{enumerate}

\subsection{Xây dựng Điều kiện Đường đi (Path Condition)}
Với mỗi đường đi $\pi$ từ đỉnh bắt đầu đến nút đích $v$, điều kiện đường đi tích lũy $PC_{\pi}$ là phép hội của tất cả các biểu thức rẽ nhánh dọc theo $\pi$:
\begin{equation}
PC_{\pi} = \bigwedge_{e = (u, w) \in \pi} \text{BranchGuard}(e)
\end{equation}
Quan hệ $R_\pi(s_0,s)$ kết hợp $PC_\pi$ với các phương trình liên hệ trạng thái đầu vào $s_0$ và trạng thái ký hiệu $s$ tại nút đích.

\section{Giai đoạn 4: Sinh Chuỗi Lý luận Có Cấu trúc từ LLM}
\label{sec:llm_structured_cot}

Để loại bỏ tính vô tổ chức của văn bản tự do, hệ thống định nghĩa một Lược đồ JSON nghiêm ngặt (Strict JSON Schema) bắt buộc mô hình LLM phải tuân thủ trong quá trình sinh chuỗi lý luận (Reasoning Trace).

\subsection{Định dạng phản hồi của LLM}
Phản hồi gồm một đoạn diễn giải bằng ngôn ngữ tự nhiên (\texttt{cot\_trace}) và danh sách các khẳng định (\texttt{claims}). Mỗi claim chứa số dòng (\texttt{line}) và biểu thức Python (\texttt{expression}).
\begin{lstlisting}[caption={Ví dụ về phản hồi JSON của LLM}]
{
  "cot_trace": "The loop counter is nonnegative.",
  "claims": [{"line": 5, "expression": "i >= 0"}]
}
\end{lstlisting}
Claim tại dòng $L$ được diễn giải trên trạng thái ngay trước khi câu lệnh tại $L$ thực thi. Ví dụ, nếu dòng $L$ gán \texttt{x = n + 3}, khẳng định về giá trị mới của \texttt{x} phải đặt ở vị trí sau phép gán. Trình thực thi ký hiệu tự xây dựng quan hệ trạng thái từ AST. Đoạn diễn giải cung cấp ngữ cảnh cho người đọc và vòng phản hồi.

\section{Giai đoạn 5: Cổng neo AST (Anchor Gate)}
\label{sec:anchor_gate}

Hàm \texttt{bind\_llm\_claims} ánh xạ số dòng trong mỗi đề xuất đến một câu lệnh AST. Claim được neo khi dòng đó định danh duy nhất một nút hợp lệ, biểu thức thuộc tập con hỗ trợ và mọi tên biến được phép sử dụng ở vị trí đó. Claim có biểu thức không phân tích cú pháp được nhận nhãn \texttt{TRANSLATION\_ERROR}. Trường hợp vị trí không xác định duy nhất, biểu thức ngoài tập con hoặc có biến ngoài phạm vi nhận nhãn \texttt{UNSUPPORTED}.

\section{Giai đoạn 6: Giao thức Hai Truy vấn Z3 (Two-Query Z3 Protocol)}
\label{sec:two_query_protocol}

Với mỗi claim đã được neo, hệ thống xét các đường đi tới nút đích. Gọi $R_\pi(s_0,s)$ là quan hệ từ trạng thái đầu vào $s_0$ đến trạng thái $s$ ngay trước câu lệnh đích trên đường $\pi$. Ký hiệu $A_\pi$ là ràng buộc tiền điều kiện mà bộ sinh nghĩa vụ đưa vào cả hai truy vấn. Hiện tại, mã đánh giá biểu thức tiền điều kiện trong môi trường ký hiệu tại đích. Để diễn giải nó là tiền điều kiện trên đầu vào, cần bảo đảm các tên được tham chiếu tại đích vẫn biểu diễn giá trị đầu vào tương ứng.

\subsection{Truy vấn 1: Kiểm tra tính khả đạt}
\begin{equation}
\Phi_{\text{reach},\pi} \equiv \exists s_0,s.\quad R_\pi(s_0,s) \land A_\pi
\end{equation}
Kết quả \texttt{UNSAT} loại đường $\pi$ khỏi bước kiểm tra tính hợp lệ. Kết quả \texttt{SAT} cho phép chạy truy vấn thứ hai trên đường này. Nếu bộ giải trả về \texttt{UNKNOWN}, hệ thống ghi nhận \texttt{UNKNOWN\_TIMEOUT}.

\subsection{Truy vấn 2: Tìm vi phạm của khẳng định}
\begin{equation}
\Phi_{\text{valid},\pi} \equiv \exists s_0,s.\quad R_\pi(s_0,s) \land A_\pi \land \neg C(s)
\end{equation}
Khi kết quả là \texttt{SAT}, hệ thống lấy mô hình để xây dựng phản ví dụ ứng viên và có thể chạy lại trên CPython. Kết quả \texttt{UNSAT} cho biết không có trạng thái vi phạm trong mô hình của đường $\pi$. Kết quả \texttt{UNKNOWN} dẫn tới \texttt{UNKNOWN\_TIMEOUT}.

\begin{figure}[H]
\centering
\begin{tikzpicture}[
    node distance=1.0cm and 1.2cm,
    block/.style={rectangle, draw=blue!80, fill=blue!5, rounded corners, thick, text centered, font=\small, minimum height=0.8cm},
    decision/.style={diamond, draw=purple!80, fill=purple!10, thick, text centered, font=\small\bfseries, inner sep=1pt},
    verdict/.style={rectangle, draw=green!60!black, fill=green!10, rounded corners, thick, text centered, font=\small\bfseries, minimum height=0.8cm},
    fail/.style={rectangle, draw=red!80, fill=red!10, rounded corners, thick, text centered, font=\small\bfseries, minimum height=0.8cm},
    arrow/.style={-{Stealth[scale=1.0]}, thick}
]
    \node[block] (claim) {Claim đã neo từ Anchor Gate};
    \node[decision, below=0.8cm of claim] (q1) {Query 1: $\Phi_{\text{reach}}$ SAT?};
    \node[fail, right=1.5cm of q1] (unreach) {Bỏ qua đường $\pi$};
    \node[decision, below=1.0cm of q1] (q2) {Query 2: $\Phi_{\text{valid}}$ SAT?};
    \node[fail, right=1.5cm of q2] (ce) {Phản ví dụ ứng viên};
    \node[verdict, below=0.8cm of q2] (ver) {Không có vi phạm trên $\pi$};

    \draw[arrow] (claim) -- (q1);
    \draw[arrow] (q1) -- node[above, font=\scriptsize] {UNSAT} (unreach);
    \draw[arrow] (q1) -- node[right, font=\scriptsize] {SAT} (q2);
    \draw[arrow] (q2) -- node[above, font=\scriptsize] {SAT (Vi phạm)} (ce);
    \draw[arrow] (q2) -- node[right, font=\scriptsize] {UNSAT (An toàn)} (ver);
\end{tikzpicture}
\caption{Hai truy vấn trên một đường đi $\pi$ khi Z3 trả kết quả SAT hoặc UNSAT}
\label{fig:two_query_flowchart}
\end{figure}

\subsection{Tổng hợp theo claim}
Nếu mọi đường được xét đều vô nghiệm ở truy vấn khả đạt, hệ thống kết luận \texttt{UNREACHABLE} khi đã bao phủ đầy đủ các đường tới nút đích. Nếu ít nhất một đường có nghiệm và truy vấn thứ hai vô nghiệm trên mọi đường có nghiệm, hệ thống có thể kết luận \texttt{VERIFIED} khi độ bao phủ đầy đủ hoặc khi có bản tóm lược vòng lặp đủ để bảo đảm độ bao phủ. Khi chỉ có một tiền tố duỗi vòng lặp hữu hạn và chưa đủ bằng chứng, hệ thống trả về \texttt{UNKNOWN\_TIMEOUT}.

\section{Giai đoạn 7: Tổng hợp Bất biến Vòng lặp và Quy nạp $k$-Induction}
\label{sec:loop_invariant_synthesis_method}

Đối với các chương trình chứa vòng lặp không thể xác định cận lặp tĩnh, việc duỗi vòng lặp sẽ dẫn đến kết quả \texttt{UNKNOWN\_TIMEOUT}. Để đạt được kết luận \texttt{VERIFIED} đầy đủ, mô-đun \texttt{LoopInvariantGenerator} kích hoạt quy trình sinh và kiểm chứng bất biến:

\begin{enumerate}
    \item \textbf{Đề xuất Ứng viên Bất biến (Candidate Proposal):} Trích xuất trạng thái trước vòng lặp, biểu thức Guard, và thân vòng lặp. Yêu cầu LLM đóng vai trò là một chuyên gia toán học, đề xuất từ 1 đến 3 biểu thức ứng viên bất biến $I_1, I_2, I_3 \in \text{QF-LIA}$. Song song với LLM, hệ thống luôn sinh sẵn các mẫu bất biến tương đối (relational) theo dạng $(x \geq 0)$, $(x \leq y)$ để bảo đảm cơ chế hoạt động khi không có khoá API.
    \item \textbf{Kiểm chứng Tính Khởi tạo:} Với mỗi ứng viên $I$, thiết lập truy vấn Z3:
    \begin{equation}
    \Phi_{\text{init}} \equiv \text{Precondition}(s_0) \land \neg I(s_0)
    \end{equation}
    Trạng thái $s_0$ được ràng buộc bằng \emph{trạng thái thực tế tại đầu vòng lặp} do bộ thực thi ký hiệu tin cậy cung cấp. Nếu $\text{Check}(\Phi_{\text{init}}) \ne \texttt{UNSAT}$, loại bỏ ứng viên $I$.
    \item \textbf{Kiểm chứng Bước Quy nạp (Inductive Step):} Thiết lập truy vấn Z3:
    \begin{equation}
    \Phi_{\text{step}} \equiv I(s) \land \text{Guard}(s) \land T_{\text{body}}(s, s') \land \neg I(s')
    \end{equation}
    Nếu $\text{Check}(\Phi_{\text{step}}) == \texttt{UNSAT}$, thì $I$ là một \textbf{Bất biến Quy nạp Hợp lệ} (Inductive Invariant).
    \item \textbf{Mở rộng $k$-Induction:} Nếu bước quy nạp thất bại ở $k=1$, hệ thống tự động tăng bậc quy nạp lên $k=2$ và $k=3$ theo công thức ở Mục \ref{sec:loop_invariants_k_induction}.
    \item \textbf{Suy luận Hậu điều kiện:} Sử dụng bất biến đã chứng minh $I$ làm tiên đề bổ sung để chứng minh các câu lệnh \texttt{assert} sau vòng lặp.
\end{enumerate}

\textbf{Phạm vi hiện thực.}
Bước quy nạp được mã hoá ký hiệu cho mọi thân vòng lặp thuộc phân đoạn số học nguyên tuyến tính (QF-LIA), bao gồm các phép gán tuyến tính tổng quát ($x \leftarrow e$ với $e$ là tổ hợp tuyến tính của các biến và hằng số, phép chia nguyên/chia lấy dư cho hằng số), các phép cập nhật gán kết hợp (\texttt{+=}, \texttt{-=}), cũng như các cấu trúc rẽ nhánh có điều kiện \texttt{if/else} lồng nhau (được hợp nhất trạng thái bằng toán tử chọn $\text{ite}(\text{cond}, s_{\text{then}}, s_{\text{else}})$ của Z3). Ngược lại, một thân vòng lặp chứa lệnh nhảy làm thay đổi luồng điều khiển sớm (\texttt{break}, \texttt{continue}, \texttt{return}), vòng lặp lồng nhau, hoặc phép toán phi tuyến giữa hai biến sẽ bị từ chối mã hoá. Trong các trường hợp đó, ứng viên bất biến bị loại bỏ và hệ thống giữ nguyên phán quyết \texttt{UNKNOWN\_TIMEOUT}. Thiết kế này bảo đảm tính \emph{đúng đắn} (soundness): chỉ khi bước quy nạp thực sự được chứng minh bằng Z3 thì bất biến mới được hợp nhất vào quan hệ trạng thái, tránh việc một mệnh đề hằng đúng (tautology) làm sai lệch phán quyết.

\textbf{Đưa bất biến vào giao thức hai truy vấn.}
Bất biến đã được chứng minh được hợp nhất (conjunct) trực tiếp vào quan hệ trạng thái $R_\pi$ của đường tới claim, trước khi dựng hai bộ giải Z3. Cụ thể:
\begin{equation}
R'_\pi = R_\pi \land I(s)
\end{equation}
Nhờ vậy, hai truy vấn \emph{reachability} và \emph{validity} suy luận trên tập trạng thái đã được tóm lược bởi bất biến thay vì trên một tiền tố giới hạn của các lần lặp, cho phép phán quyết \texttt{VERIFIED} thay cho \texttt{UNKNOWN\_TIMEOUT} đối với các claim nằm trong hoặc sau vòng lặp.

\section{Giai đoạn 8: Replay Phản ví dụ trong Sandbox CPython \& Vòng Phản hồi Hồi cứu}
\label{sec:replay_and_retrospective_feedback}

Khi Z3 phát ra phán quyết \texttt{COUNTEREXAMPLE}, hệ thống không lập tức dừng lại mà kích hoạt tầng kiểm tra chéo thực nghiệm và chu trình tự sửa sai đóng kín.

\subsection{Cơ chế Replay trong Sandbox CPython}
Thuật toán trích xuất các giá trị gán của biến đầu vào từ mô hình Z3 $\mathcal{M}$:
\begin{equation}
\mathbf{x}_{\text{concrete}} = \{ v \mapsto \mathcal{M}.\text{eval}(v_{\text{initial}}).\text{as\_long}() \}
\end{equation}

Hệ thống sinh một kịch bản kiểm thử độc lập (test harness) và khởi chạy trong một tiến trình con CPython biệt lập với môi trường runtime:
\begin{itemize}
    \item \textbf{Replay thành công (\texttt{replay = True}):} Tiến trình CPython tới vị trí neo và xác nhận biểu thức claim sai tại đó. Với claim trích từ một câu lệnh \texttt{assert}, điều này tương ứng với vi phạm assertion tại vị trí đó.
    \item \textbf{Replay thất bại (\texttt{replay = False}):} Lượt chạy không xác nhận được vi phạm, chẳng hạn vì không đến vị trí neo hoặc hết thời gian. Hệ thống hạ xuống \texttt{UNKNOWN\_TIMEOUT} nếu có đường đi chưa được bao phủ đầy đủ. Trong trường hợp còn lại, hệ thống giữ \texttt{COUNTEREXAMPLE} kèm cờ replay thất bại.
\end{itemize}

\subsection{Vòng Phản hồi Hồi cứu (Retrospective Feedback Loop)}
Khi bật cơ chế phản hồi và CPython xác nhận vi phạm tại vị trí neo, tùy chọn \texttt{capture\_state} của mô-đun replay cho phép ghi nhận trạng thái các biến để tạo ngữ cảnh chẩn đoán.
\begin{equation}
\mathcal{C}_{\text{feedback}} = \langle \text{Input: } \mathbf{x}_{\text{concrete}}, \quad \text{Crash\_Line: } L, \quad \text{State\_at\_Crash: } s_{\text{concrete}}, \quad \text{Failed\_Claim: } C \rangle
\end{equation}

Quy trình tự sửa sai diễn ra trong tối đa 3 vòng lặp (Rounds 1--3):
\begin{enumerate}
    \item Gửi $\mathcal{C}_{\text{feedback}}$ ngược lại cho LLM với chỉ thị: \textit{``Phản ví dụ sau đã được CPython xác nhận gây sụp đổ chương trình với đầu vào $\mathbf{x}_{\text{concrete}}$. Hãy phân tích nguyên nhân vì sao khẳng định ban đầu bị sai, và đề xuất lại khẳng định hoặc bất biến đã được sửa đổi.''}
    \item Nhận phản hồi mới từ LLM, đưa qua Anchor Gate để kiểm định lại.
    \item Chạy lại Giao thức Hai Truy vấn Z3 trên claim mới.
    \item Nếu claim mới vượt qua bước kiểm tra lại trên các đường đi đã mô hình hóa, kết quả \texttt{VERIFIED} được gắn với nội dung claim mới. Phản ví dụ tiếp tục áp dụng cho claim ban đầu. Nếu không chấp nhận được đề xuất mới, hệ thống giữ kết quả của claim ban đầu.
\end{enumerate}

\section{Cơ sở Tính toán Tin cậy (TCB) và Phân tích Tính Đúng đắn (Soundness)}
\label{sec:tcb_soundness}

\begin{definition}[Trusted Computing Base -- TCB]
Cơ sở tính toán tin cậy của hệ thống kiểm chứng bao gồm: (1) Trình phân tích cú pháp Python chuẩn (\texttt{ast}); (2) Bộ lọc Safe Subset; (3) Bộ thực thi ký hiệu SSA; (4) Cổng neo Anchor Gate; (5) Trình giải Z3 SMT; và (6) Trình thông dịch CPython.
\end{definition}

\begin{theorem}[Tính đúng đắn tương đối của nghĩa vụ kiểm chứng]
Giả sử công thức trạng thái bao phủ mọi lần thực thi cụ thể đạt đến vị trí neo dưới tiền điều kiện đã công bố. Giả sử $C(s)$ biểu diễn đúng khẳng định tại vị trí đó và bộ giải trả kết quả đúng. Nếu truy vấn tìm vi phạm vô nghiệm trên mọi đường đi có nghiệm, thì khẳng định đúng tại mọi lần thực thi cụ thể đạt đến vị trí neo.
\end{theorem}

\begin{proof}
Giả sử có một lần thực thi đạt đến vị trí neo nhưng vi phạm $C$. Theo giả thiết bao phủ, lần thực thi này cho một phép gán thỏa mãn $R_\pi \land A_\pi$ đối với ít nhất một đường đi $\pi$. Vì lần thực thi vi phạm $C$, phép gán đó cũng thỏa mãn $\neg C(s)$. Khi đó truy vấn tìm vi phạm trên $\pi$ có nghiệm, mâu thuẫn với giả thiết.
\end{proof}

Lập luận này áp dụng cho claim được kiểm tra và phụ thuộc vào độ bao phủ của mô hình, cách diễn giải tiền điều kiện và độ chính xác của phép dịch. Với đường đi được xấp xỉ, nghiệm của công thức khả đạt chỉ chứng minh khả đạt trong mô hình.
